\documentclass[pdf, unicode, 12pt, a4paper, oneside, fleqn]{article}

% оформление

\usepackage[utf8]{inputenc}
\usepackage[T2A]{fontenc}
\usepackage[english,russian]{babel}

\usepackage{geometry}
\geometry{
  a4paper,
  left=2.54cm,
  right=2.54cm,
  top=2.54cm,
  bottom=2.54cm,
  headheight=0mm,
  headsep=0mm,
  footnotesep=0mm,
  footskip=0mm,
}

\usepackage{nopageno}
\usepackage{tempora}
\renewcommand{\ttdefault}{cmtt}

\usepackage{titlesec}
\titleformat*{\section}{\large\bfseries}
\titlespacing\section{0pt}{12pt plus 4pt minus 2pt}{0pt plus 2pt minus 2pt}
\setcounter{secnumdepth}{0}

\usepackage{enumitem}
\frenchspacing

\usepackage{amsmath}
\usepackage{amssymb}
\usepackage{hyperref}
\usepackage{longtable}
\usepackage[table]{xcolor}
\usepackage{array}
\usepackage{color}
\usepackage{xcolor}
\usepackage{csquotes}

\DeclareQuoteStyle{russian}
	{\guillemotleft}{\guillemotright}[0.025em]
	{\quotedblbase}{\textquotedblleft}
\ExecuteQuoteOptions{style=russian}

\usepackage{graphicx}
\usepackage{listings}
\usepackage{alltt}

\newcommand{\MYhref}[3][blue]{\href{#2}{\color{#1}{#3}}}

\lstset{
  tabsize=2,
  breaklines=true,
  columns=fullflexible,
  flexiblecolumns=true,
  numbers=left,
  numberstyle={\footnotesize},
  extendedchars=true,
  inputencoding=utf8,
  frame=single,
  basicstyle=\footnotesize\ttfamily,
  keywordstyle=\bfseries\color{green!40!black},
  commentstyle=\itshape\color{gray},
  identifierstyle=\color{black},
  stringstyle=\color{blue},
  showstringspaces=false
}

\def\@xobeysp{ }
\def\verbatim@processline{\hspace{1.2cm}\raggedright\the\verbatim@line\par}

\parindent=0pt
\parskip=3pt

\definecolor{lightgray}{gray}{0.9}
\renewcommand{\thesubsection}{\arabic{subsection}}


\begin{document}

\graphicspath{ {./img/} }


% титульник

\begin{titlepage}
\begin{center}
\large
МОСКОВСКИЙ АВИАЦИОННЫЙ ИНСТИТУТ (НАЦИОНАЛЬНЫЙ ИССЛЕДОВАТЕЛЬСКИЙ УНИВЕРСИТЕТ)

\vspace{20pt}

Институт №8 <<Компьютерные науки и прикладная математика>>

Кафедра 806 <<Вычислительная математика и программирование>>
\end{center}

\vspace{60pt}

\begin{center}
\bfseries
\large
Отчет по лабораторной работе \textnumero 2 по курсу \enquote{Операционные системы}
\end{center}

\vspace{40pt}

\begin{center}
\bfseries
\large
Тема: \enquote{Управление процессами и межпроцессное взаимодействие посредством каналов (Pipe)}
\end{center}

\vfill

\begin{flushright}
\large
\begin{tabular}{rl}
Выполнил(а): & А.\,Е. Пронина \\
Группа: & М8О-203БВ-25 \\
Преподаватель: & Е.\,С. Миронов \\
\end{tabular}
\end{flushright}

\vspace{40pt}

\begin{center}
\large
Москва, \the\year
\end{center}
\end{titlepage}

\pagebreak


\section{Условие}

{\bfseries Цель работы:} приобретение практических навыков управления процессами в операционной системе и обеспечения обмена данными между ними посредством неименованных каналов (pipe), а также перенаправления стандартных потоков ввода-вывода дочерних процессов.

{\bfseries Задание:} родительский процесс создает два дочерних процесса. Первой и второй строкой пользователь в консоль родительского процесса вводит имена файлов, которые будут использованы для открытия на запись для child1 и child2 соответственно. Родительский и дочерний процессы представлены разными программами. Родительский процесс принимает от пользователя строки произвольной длины и пересылает их в pipe1 или в pipe2 в зависимости от правила фильтрации. Процессы child1 и child2 производят работу над строками и пишут результаты своей работы в стандартный вывод (перенаправленный в файлы).

{\bfseries Вариант:} 18 (Группа 5). Правило фильтрации: нечетные строки отправляются в pipe1, четные в pipe2. Дочерние процессы удаляют все гласные из строк.

\pagebreak

\section{Метод решения}
Программа реализована с использованием кроссплатформенного подхода, абстрагирующего системные вызовы POSIX (Linux/macOS) и Win32 API (Windows). 

Алгоритм работы родительского процесса:
\begin{enumerate}
    \item Запрос имен двух выходных файлов у пользователя и их открытие на запись с флагами создания/усечения (\texttt{O\_WRONLY | O\_CREAT | O\_TRUNC} / \texttt{CREATE\_ALWAYS}).
    \item Создание двух неименованных каналов (\texttt{pipe} / \texttt{CreatePipe}).
    \item Двукратный вызов функции создания дочернего процесса. Перед запуском исполняемого файла дочернего процесса стандартный поток ввода (\texttt{stdin}) перенаправляется на читающий конец соответствующего канала, а стандартный поток вывода (\texttt{stdout}) — на дескриптор открытого файла. В POSIX это достигается через \texttt{dup2}, в Windows — через структуру \texttt{STARTUPINFO} с флагом \texttt{STARTF\_USESTDHANDLES}.
    \item Родительский процесс закрывает все ненужные ему концы каналов (читающие концы обоих каналов), чтобы избежать утечек дескрипторов.
    \item В цикле считываются строки из \texttt{stdin}. Нечетные строки записываются в \texttt{pipe1}, четные — в \texttt{pipe2}.
    \item После получения признака конца файла (EOF) от пользователя, родительский процесс закрывает пишущие концы каналов. Это критически важно: только когда все дескрипторы пишущего конца закрыты, системный вызов \texttt{read} в дочернем процессе вернет 0 (EOF), позволяя ему корректно завершить цикл чтения.
    \item Ожидание завершения дочерних процессов (\texttt{waitpid} / \texttt{WaitForSingleObject}) и вывод их кодов возврата.
\end{enumerate}

\section{Описание программы}
Проект разделен на модули для обеспечения чистоты архитектуры и соответствия инженерным требованиям:
\begin{itemize}
    \item \texttt{headers/*.h}: объявление кроссплатформенных структур и прототипов функций.
    \item \texttt{realization/*\_posix.c} и \texttt{realization/*\_win.c}: реализация системно-зависимых функций (\texttt{open\_file\_for\_write}, \texttt{create\_pipe}, \texttt{create\_child\_process}).
    \item \texttt{parent.c}: основная логика оркестрации процессов.
    \item \texttt{child.c}: логика обработки данных (удаление гласных).
\end{itemize}

Ключевой фрагмент перенаправления потоков в дочернем процессе (POSIX):
\begin{lstlisting}[language=C]
if (pid == 0) { // Дочерний процесс
    if (stdin_handle != -1) {
        if (dup2((int)stdin_handle, STDIN_FILENO) == -1) {
            _exit(1);
        }
    }
    if (stdout_handle != -1) {
        if (dup2((int)stdout_handle, STDOUT_FILENO) == -1) {
            _exit(1);
        }
    }
    execvp(executable, argv);
    _exit(1);
}
\end{lstlisting}

\pagebreak


\section{Результаты работы}

Программа была протестирована на наборе из шести строк, содержащих как гласные, так и согласные буквы. Ввод строк завершался нажатием Ctrl+D (сигнал EOF). Ниже приведены примеры работы системы.

\subsection*{Консольный вывод родительского процесса}
В консоли отображаются только сообщения родительского процесса: его PID, запрос имен файлов, PID созданных дочерних процессов и коды их завершения после окончания работы.

\begin{lstlisting}[language=bash, caption=Вывод в консоль родительского процесса]
$ ./parent
[Parent] PID: 89233
Enter output file for Child1: ch1
Enter output file for Child2: ch2
[Parent] Child1 created with PID: 89265
[Parent] Child2 created with PID: 89266
Enter lines (Ctrl+D / Ctrl+Z to finish):
fdg
sdfon
ainaaon
aaass
aiaoaiau
jnv
[Parent] Child1 (PID: 89265) exit code: 0
[Parent] Child2 (PID: 89266) exit code: 0
[Parent] Results written to: ch1, ch2
\end{lstlisting}

\subsection*{Содержимое выходных файлов}

Поскольку стандартный поток вывода (\texttt{stdout}) дочерних процессов перенаправлен в файлы, служебные сообщения о запуске и завершении, а также результаты обработки строк, записываются непосредственно в файлы.

Файл \texttt{ch1} содержит результаты обработки \textbf{нечетных} строк (1, 3, 5):
\begin{itemize}
    \item Строка 1: \texttt{fdg} $\rightarrow$ \texttt{fdg} (гласных нет)
    \item Строка 3: \texttt{ainaaon} $\rightarrow$ \texttt{nn} (удалены a, i, a, a, o)
    \item Строка 5: \texttt{aiaoaiau} $\rightarrow$ \texttt{""} (все символы — гласные, остается только символ перевода строки)
\end{itemize}

\begin{lstlisting}[language=text, caption=Содержимое файла ch1]
[Child1] PID: 89265, started
fdg
nn

[Child1] PID: 89265, finished
\end{lstlisting}

Файл \texttt{ch2} содержит результаты обработки \textbf{четных} строк (2, 4, 6):
\begin{itemize}
    \item Строка 2: \texttt{sdfon} $\rightarrow$ \texttt{sdfn} (удалена o)
    \item Строка 4: \texttt{aaass} $\rightarrow$ \texttt{ss} (удалены a, a, a)
    \item Строка 6: \texttt{jnv} $\rightarrow$ \texttt{jnv} (гласных нет)
\end{itemize}

\begin{lstlisting}[language=text, caption=Содержимое файла ch2]
[Child2] PID: 89266, started
sdfn
ss
jnv
[Child2] PID: 89266, finished
\end{lstlisting}

\subsection*{Анализ результатов}
Результаты полностью соответствуют заданию Варианта 18:
\begin{enumerate}
    \item Нечетные строки корректно перенаправлены в \texttt{pipe1} и обработаны дочерним процессом 1, результат записан в файл \texttt{ch1}.
    \item Четные строки корректно перенаправлены в \texttt{pipe2} и обработаны дочерним процессом 2, результат записан в файл \texttt{ch2}.
    \item Все гласные буквы (a, e, i, o, u) удалены из строк.
    \item Служебные сообщения дочерних процессов (PID, статус запуска/завершения) присутствуют в выходных файлах, что подтверждает корректность перенаправления потока \texttt{stdout} в файлы посредством системных вызовов \texttt{dup2} (POSIX) / \texttt{STARTUPINFO} (Win32).
    \item Оба дочерних процесса завершились с кодом возврата 0, что свидетельствует об отсутствии ошибок выполнения и корректном получении EOF после закрытия родительским процессом пишущих концов каналов.
\end{enumerate}
\pagebreak


\section*{Диагностика системных вызовов (strace)}
Для подтверждения использования нативных системных примитивов ОС был проведен анализ профиля системных вызовов программы с помощью утилиты \texttt{strace}. Запуск осуществлялся с флагом \texttt{-f} для трассировки всех дочерних процессов и фильтрацией по категориям \texttt{process, ipc, write}.

Ниже приведены ключевые фрагменты реального лога выполнения, демонстрирующие жизненный цикл процессов и межпроцессное взаимодействие:

\begin{lstlisting}[language=bash, caption=Фрагмент вывода strace (ключевые системные вызовы)]
strace -f -e trace=process,ipc,write  ./parent
execve("./parent", ["./parent"], 0x7ffdfdb7dcc8 /* 38 vars */) = 0
write(1, "[Parent] PID: 91777\n", 20[Parent] PID: 91777
)   = 20
write(1, "Enter output file for Child1: ", 30Enter output file for Child1: ) = 30
child1
write(1, "Enter output file for Child2: ", 30Enter output file for Child2: ) = 30
child2
clone(child_stack=NULL, flags=CLONE_CHILD_CLEARTID|CLONE_CHILD_SETTID|SIGCHLDstrace: Process 91850 attached
, child_tidptr=0x7f572d39fa10) = 91850
[pid 91777] clone(child_stack=NULL, flags=CLONE_CHILD_CLEARTID|CLONE_CHILD_SETTID|SIGCHLDstrace: Process 91851 attached
 <unfinished ...>
[pid 91850] execve("./child", ["./child", "child", "1"], 0x7ffddee99fe8 /* 38 vars */ <unfinished ...>
[pid 91777] <... clone resumed>, child_tidptr=0x7f572d39fa10) = 91851
[pid 91777] write(1, "[Parent] Child1 created with PID"..., 40[Parent] Child1 created with PID: 91850
) = 40
[pid 91777] write(1, "[Parent] Child2 created with PID"..., 40[Parent] Child2 created with PID: 91851
) = 40
[pid 91851] execve("./child", ["./child", "child", "2"], 0x7ffddee99fe8 /* 38 vars */ <unfinished ...>
[pid 91777] write(1, "Enter lines (Ctrl+D / Ctrl+Z to "..., 41Enter lines (Ctrl+D / Ctrl+Z to finish):
) = 41
[pid 91851] write(1, "[Child2] PID: 91851, started\n", 29 <unfinished ...>
[pid 91850] write(1, "[Child1] PID: 91850, started\n", 29 <unfinished ...>

aaaa
[pid 91777] write(4, "aaaa\n", 5)       = 5
[pid 91850] write(1, "\n", 1)           = 1

ooooo
[pid 91777] write(6, "ooooo\n", 6)      = 6
[pid 91851] write(1, "\n", 1)           = 1


lalalal
[pid 91777] write(4, "lalalal\n", 8)    = 8
[pid 91850] write(1, "llll\n", 5)       = 5
aaiaaoa
[pid 91777] write(6, "aaiaaoa\n", 8)    = 8
[pid 91851] write(1, "\n", 1)           = 1
[pid 91777] wait4(91850 <unfinished ...>
[pid 91850] write(1, "[Child1] PID: 91850, finished\n", 30 <unfinished ...>
[pid 91851] write(1, "[Child2] PID: 91851, finished\n", 30) = 30
[pid 91851] exit_group(0 <unfinished ...>
[pid 91850] exit_group(0 <unfinished ...>
[pid 91851] <... exit_group resumed>)   = ?
[pid 91850] <... exit_group resumed>)   = ?
[pid 91850] +++ exited with 0 +++
[pid 91777] <... wait4 resumed>, [{WIFEXITED(s) && WEXITSTATUS(s) == 0}], 0, NULL) = 91850
[pid 91851] +++ exited with 0 +++
--- SIGCHLD {si_signo=SIGCHLD, si_code=CLD_EXITED, si_pid=91850, si_uid=1000, si_status=0, si_utime=0, si_stime=0} ---
wait4(91851, [{WIFEXITED(s) && WEXITSTATUS(s) == 0}], 0, NULL) = 91851
write(1, "[Parent] Child1 (PID: 91850) exi"..., 42[Parent] Child1 (PID: 91850) exit code: 0
) = 42
write(1, "[Parent] Child2 (PID: 91851) exi"..., 42[Parent] Child2 (PID: 91851) exit code: 0
) = 42
write(1, "[Parent] Results written to: chi"..., 44[Parent] Results written to: child1, child2
) = 44
exit_group(0)                           = ?
+++ exited with 0 +++
\end{lstlisting}

Данный лог наглядно подтверждает, что программа использует требуемые низкоуровневые системные вызовы ядра Linux (\texttt{clone}, \texttt{execve}, \texttt{write}, \texttt{wait4}) для управления процессами и организации IPC, без применения высокоуровневых библиотечных оберток.

\section{Выводы}

В ходе выполнения лабораторной работы была спроектирована и реализована кроссплатформенная система межпроцессного взаимодействия (IPC) с использованием неименованных каналов. Практическая реализация позволила глубже изучить механизмы, лежащие в основе управления процессами и обмена данными в операционных системах.

Ключевым теоретическим аспектом, подтвержденным на практике, является природа каналов (\texttt{pipe}) как однонаправленных потоков байтов, буферизуемых в адресном пространстве ядра. Системный вызов \texttt{pipe()} (POSIX) или \texttt{CreatePipe()} (Win32) создает пару файловых дескрипторов: один для чтения, другой для записи. Данные, записанные в один конец канала, могут быть прочитаны из другого в порядке FIFO. Важной особенностью является то, что системный вызов \texttt{read()} блокирует процесс-читатель до тех пор, пока не появятся данные или пока все дескрипторы пишущего конца не будут закрыты. Именно поэтому критически важно, чтобы родительский процесс закрыл \textit{все} копии дескрипторов, связанных с пишущим концом канала после \texttt{fork()}, иначе дочерний процесс будет бесконечно блокироваться в ожидании данных, что приведет к взаимоблокировке (deadlock).

Механизм создания дочерних процессов в POSIX реализуется парой вызовов \texttt{fork()} и \texttt{exec()}. Вызов \texttt{fork()} создает точную копию текущего процесса, наследующую все открытые файловые дескрипторы, что является основой для последующего перенаправления потоков. Вызов \texttt{execvp()} заменяет образ памяти дочернего процесса на исполняемый файл, сохраняя при этом перенаправленные дескрипторы. В Windows аналогичная функциональность инкапсулирована в едином системном вызове \texttt{CreateProcess()}, который принимает структуру \texttt{STARTUPINFO} с флагом \texttt{STARTF\_USESTDHANDLES} для явного указания дескрипторов стандартных потоков.

Перенаправление стандартных потоков ввода-вывода достигается посредством системного вызова \texttt{dup2()} (POSIX), который дублирует файловый дескриптор, заменяя им стандартные дескрипторы \texttt{STDIN\_FILENO} (0) и \texttt{STDOUT\_FILENO} (1). Этот механизм демонстрирует принцип прозрачности файловых дескрипторов в UNIX-подобных системах: дочерняя программа (\texttt{child.c}) работает исключительно с абстракциями \texttt{stdin} и \texttt{stdout}, не содержа в себе никакой логики работы с каналами или файлами. Операционная система сама маршрутизирует данные между процессами на уровне ядра.

Ожидание завершения дочерних процессов реализовано через \texttt{waitpid()} (POSIX) и \texttt{WaitForSingleObject()} (Win32). Эти вызовы не только блокируют родительский процесс до завершения дочернего, но и позволяют получить код возврата, а также освободить запись в таблице процессов ядра. Без вызова \texttt{wait()} завершенный дочерний процесс остается в состоянии "зомби" (zombie), занимая место в таблице процессов до завершения родителя.

Также была отработана инженерная практика безопасного управления ресурсами: каждый системный вызов (\texttt{fork}, \texttt{pipe}, \texttt{open}) сопровождается проверкой кода возврата, а выделенные дескрипторы гарантированно освобождаются при любом сценарии выполнения, включая аварийное завершение. Это предотвращает утечки ресурсов и появление зомби-процессов, что соответствует требованиям к надежному системному программированию.
\pagebreak


\section{Список литературы}
\begin{enumerate}
    \item Таненбаум Э., Бос Х. Современные операционные системы. – 4-е изд. – СПб.: Питер, 2015.
\end{enumerate}

\end{document}